% Generated draft: requires mathematical and source review.
\documentclass[11pt]{article}
\usepackage[margin=25mm]{geometry}
\usepackage{fontspec}
\setmainfont{Latin Modern Roman}
\usepackage{amsmath,amssymb,mathtools,unicode-math}
\setmathfont{Latin Modern Math}
\usepackage{enumitem,needspace,xurl,hyperref}
\setlength{\parindent}{0pt}
\setlength{\parskip}{5pt}
\setlength{\emergencystretch}{3em}
\newcommand{\N}{\mathbb{N}}
\newcommand{\Z}{\mathbb{Z}}
\newcommand{\Q}{\mathbb{Q}}
\newcommand{\R}{\mathbb{R}}
\newcommand{\C}{\mathbb{C}}
\newcommand{\F}{\mathbb{F}}
\newcommand{\A}{\mathbb{A}}
\newcommand{\PP}{\mathbb{P}}
\newcommand{\E}{\mathbb{E}}
\newcommand{\eps}{\varepsilon}
\newcommand{\dd}{\,\mathrm{d}}
\begin{document}
\section*{UnsolvedMath ID: 20001904}
\textbf{Source problem code:} catalogue entry
\subsection{Definitions and mathematical statement}
The source asks whether \(d(\alpha(t),\beta(t))\to0\) as \(t\to-\infty\) implies \(d(\alpha(t),\beta(t))\to\infty\) as \(t\to\infty\). It omits preceding hypotheses on the space and geodesics. Distinctness is essential, since \(\alpha=\beta\) is an immediate exception; the distance must also be in the intended noncompact covering space, not a compact base of bounded diameter. The aid's precise positive result concerns distinct unit-speed geodesic lines with synchronized parameters.

A geodesic space is Busemann-convex if the distance between any two affinely parametrized geodesic segments is convex in the common parameter. For unit-speed lines \(\alpha,\beta:\R\to X\), put \(f(t)=d(\alpha(t),\beta(t))\) and suppose \(\lim_{t\to-\infty}f(t)=0\). The aid reports that \(f\) is nondecreasing and either identically zero or tends to infinity. More precisely, any \(a<b\) with
\[
 m=\frac{f(b)-f(a)}{b-a}>0
\]
gives, for every \(t\ge b\),
\[
 f(b)+m(t-b)\le f(t)\le f(b)+2(t-b).
\]
The upper slope two uses unit speed for both lines and is reported universally sharp. A nonidentical pair has some positive secant under the backward-limit condition, giving linear forward escape. Coincidence on a past interval is not excluded; the theorem permits branching spaces.

This answers the reconstructed question in complete simply connected Riemannian manifolds with sectional curvature \(K\le0\), whose line-distance functions have the required convexity. Under the weaker complete simply connected no-focal-points hypothesis, the aid reports only that a distinct such pair cannot remain bounded forward: boundedness in both directions would give a flat strip and contradict the zero backward separation. This gives unboundedness, not by itself the limit \(+\infty\) or a positive linear escape slope. The full forward-limit question under no conjugate points alone remains reported unresolved by the aid. The absent original hypotheses should not be recovered by silently substituting curvature \(K\le0\).

The source motivates the questions by a hoped-for sphere-at-infinity theory. The appended Section 7.2 recalls the flat-strip theorem for parallel geodesics in simply connected nonpositive curvature, its extension without focal points, and its failure without conjugate points, citing [Burns1992] and an unpublished Kleiner example. It says Kleiner perturbs Heintze's example, Example 4 of [Ba-Br-Eb1985], with \(\Z^2\) in the fundamental group but no corresponding flat torus. These are flat-strip examples, not supplied counterexamples to the strong-asymptotic question. The reference to Proposition 4 of [Eschenburg1977] concerns a following parallel-postulate question whose text is not included here; its extra hypotheses are not invented or applied to the present problem.
\subsection{Short English statement}
The aid proves forward divergence for distinct backward-asymptotic unit-speed lines under convex distance, with an explicit secant lower bound and sharp upper slope two. Without focal points it reports only forward unboundedness; the weaker no-conjugate-points limit question remains unresolved in its account.
\subsection{Sources}
\begin{itemize}
\item Local UnsolvedMath registry, permanent ID \(20001904\), source code \texttt{catalogue entry}; \texttt{aim-geometry-notes.json}, \texttt{notes[241]}, tag \texttt{question}.
\item AIM workshop ``Geodesics,'' recorded item 7.1, extracted Question 7.1.2: \url{https://aimath.org/pastworkshops/geodesicsproblems.pdf}. Source type PDF; statement status \texttt{exact}. The extraction includes a page-16 header naming Keith Burns and Vladimir S. Matveev and the following Section 7.2. Only the supplied text was used; no PDF was accessed or compiled.
\item Source labels [Burns1992], [Ba-Br-Eb1985], and [Eschenburg1977], Proposition 4; Kleiner unpublished example; Heintze example. No authoritative full bibliographic entries, exact missing context, or complete proof artifacts are supplied. Original remarks and literature fields are empty.
\item Research attempt 1: \texttt{partial\_result}, status at run \texttt{open}, rigor label \texttt{proved}, novelty confidence low; registry status \texttt{partially\_solved}, classification \texttt{PARTIAL-PROGRESS}. Created and updated 14 August 2026. Machine-generated aid, not peer reviewed; independently verify before citation.
\item Ulam AI / UnsolvedMath Contributors; CC BY 4.0.
\end{itemize}
\end{document}
